%------------------------------------------------------------------------------%
\documentclass[fleqn,utf8,aspectratio=169,12pt]{beamer}

\usepackage{gaal}

\title{Definição de Vetores}

%------------------------------------------------------------------------------%
\begin{document}

\addtitle

%------------------------------------------------------------------------------%
\section{Vetores}

%------------------------------------------------------------------------------%
\begin{frame}{Vetor}
  Vetores são ``setinhas''

  \pause
  \bigskip
  Precisamos dizer isso de forma \emph{precisa}
  \pause
  e que nos permita fazer \emph{cálculos}
\end{frame}

\framedgraphic{Mapa de Ventos}{mapa_ventos}{}
\framedgraphic{Mapa de Ventos}{ventos_regiao_sul}{}
\framedgraphic{Fluxo em um Aerofólio}{aerofolio}{}
\framedgraphic{Linhas de Fluxo}{aviao}{}

%------------------------------------------------------------------------------%
\section{Definição Geométrica}

%------------------------------------------------------------------------------%
\begin{frame}{Definição geométrica de vetor}
  \onslide<+->{Um vetor é representado por um segmento de reta orientado}

  \onslide<+->{Representamos um vetor por uma seta
  \[
    \overrightarrow{AB}
  \]}
  \onslide<+->{O ponto $A$ é a origem do vetor}

  \onslide<+->{O ponto $B$ é sua extremidade}

  \onslide<+->{O vetor $\overrightarrow{AB}$ representa um \emph{deslocamento}}
  \onslide<+->{do ponto $A$ para o ponto $B$}
\end{frame}

%------------------------------------------------------------------------------%
\begin{frame}{Características dos vetores}
  \onslide<+->{Um vetor possui três características fundamentais}
  \begin{description}[<+->]
      \item[módulo] seu comprimento
      \item[direção] a direção da reta que o contém
      \item[sentido] indicado pela orientação da seta
  \end{description}
\end{frame}

%------------------------------------------------------------------------------%
%\begin{frame}{Representação geométrica}
%\begin{center}
%\begin{tikzpicture}[scale=1.2,>=Stealth]
%    \draw[->] (0,0) -- (5,0) node[right] {$x$};
%    \draw[->] (0,0) -- (0,3) node[above] {$y$};
%
%    \draw[->,thick] (1,0.7) -- (4,2.3);
%
%    \fill (1,0.7) circle (2pt) node[below left] {$A$};
%    \fill (4,2.3) circle (2pt) node[above right] {$B$};
%
%    \node at (2.7,1.8) {$\overrightarrow{AB}$};
%\end{tikzpicture}
%\end{center}
%\end{frame}

%------------------------------------------------------------------------------%
\begin{frame}{Vetores iguais}
  \onslide<+->{Dois vetores são considerados iguais quando possuem}
  \begin{itemize}[<+->]
      \item o mesmo módulo
      \item a mesma direção
      \item o mesmo sentido
  \end{itemize}

  \onslide<+->{A posição em que o vetor é desenhado \emph{não} determina o vetor}

  \onslide<+->{Podemos deslocar um vetor paralelamente sem alterar suas características}
\end{frame}

%------------------------------------------------------------------------------%
\begin{frame}{Vetor representado a partir da origem}
  Frequentemente, representamos um vetor partindo da origem
  \pause
  \[
    O = (0, 0)
  \]
  \pause
  Notação
  \[
    \vec{v} = \overrightarrow{OP}
  \]
  \pause
  Se
  \[
    P = ( 3, 2 )
  \]
  \pause
  o vetor pode ser representado por
  \[
    \vec{v}
    = ( 3, 2 )
    = \vetor{3}{2}
  \]
  \pause
  Essa representação é conveniente para realizar cálculos
\end{frame}

%------------------------------------------------------------------------------%
\begin{frame}{Independência do Ponto de Origem}
  O vetor não está associado exclusivamente à origem

  \pause
  O mesmo vetor pode ser representado a partir de qualquer ponto do plano
\end{frame}

%------------------------------------------------------------------------------%
%------------------------------------------------------------------------------%
\begin{question}
  Represente o vetor
  \(
    \vec{v} = ( 2, 3 )
  \)
  partindo de
  \(
    A = ( 2, -1 )
  \)
\end{question}

%------------------------------------------------------------------------------%
\begin{resolution}{Solução}
\begin{columns}
\column{0.6\linewidth}
  O vetor \(\vec{v} = ( 2, 3 )\) define

  \pause
  \qquad deslocamento de 2 unidades para a direita

  \pause
  \qquad deslocamento de 3 unidades para cima

  \pause
  \bigskip
  Partindo de \(A = (2, -1)\)

  \pause
  \bigskip
  Avançamos para \(( 4, -1)\)

  \pause
    \bigskip
  Depois para \(( 4, 2 )\)
\column{0.4\linewidth}
  \pause
  \includegraphics{E-01-Vetores-ex-01}
\end{columns}
\end{resolution}

%------------------------------------------------------------------------------%
\endinput


%------------------------------------------------------------------------------%
\section{Definição algébrica}

%------------------------------------------------------------------------------%
\begin{frame}{Definição Algébrica}
  No plano, $\R^2$, cada vetor pode ser representado por um par ordenado
  \[
    \vec{v}
    = ( v_1, v_2 )
    = \vetor{v_1}{v_2}
  \]
  \pause
  No espaço tridimensional, $\R^3$,
  \[
    \vec{v}
    = ( v_1, v_2, v_3 )
    = \vetortri{v_1}{v_2}{v_3}
  \]
\end{frame}

%------------------------------------------------------------------------------%
\begin{frame}{Definição Algébrica}
  Um vetor de $\R^n$ é uma $n$-upla
  \[
    \vec{v}
    = ( v_1, v_2, \ldots, v_n )
    = \begin{pmatrix}[c]
        v_1 \\
        v_2 \\
        \vdots \\
        v_n
    \end{pmatrix}
  \]
  \pause
  Os números $v_1$, $v_2$, \dots, $v_n$ são as \emph{componentes} do vetor
\end{frame}

%------------------------------------------------------------------------------%
%------------------------------------------------------------------------------%
\begin{question}
  Encontre o vetor $\vec{v} = \overrightarrow{AB}$ determinado pelos pontos
  \[
    A = ( 2, 3 )
  \]
  \[
    B = ( 1, 4 )
  \]
\end{question}

%------------------------------------------------------------------------------%
\begin{resolution}{Solução}
  O deslocamento na coordenada $x$ entre os pontos $A$ e $B$ é
  \pause
  \[
    v_x
    \pause
    \;=\; 1 - 2
    \pause
    \;=\; -1
  \]
  \pause
  O deslocamento na coordenada $y$ entre os pontos $A$ e $B$ é
  \[
    \pause
    v_y
    \pause
    \;=\; 4 - 3
    \pause
    \;=\; 1
  \]
  \pause
  Portanto
  \[
    \pause
    \vec{v}
    \pause
    \;=\; \vetor{v_x}{v_y}
    \pause
    \;=\; \vetor{-1}{1}
  \]
\end{resolution}

%------------------------------------------------------------------------------%
\endinput

%------------------------------------------------------------------------------%
\begin{question}
  Dado o vetor $\vec{u} = \overrightarrow{PQ}$ com
  \(
    P = ( 2, -1 )
  \)
  q
  \(
    Q = ( 5, 3 )
  \)
  determine

  1) o vetor $\vec{u}$

  2) o ponto na extremidade do vetor $\vec{u}$ partindo da origem

  3) o ponto na extremidade do vetor $\vec{u}$ partindo do ponto $( -1, -2)$
\end{question}

%------------------------------------------------------------------------------%
\begin{resolution}{Solução}
  \[
    \vec{u}
    \pause
    \;=\; \overrightarrow{AB}
    \pause
    \;=\; \vetor{5 - 2}{3 - (-1)}
    \pause
    \;=\; \vetor{3}{4}
  \]

  \pause
  O vetor $\vec{u}$ partindo da origem termina no ponto $(3, 4)$

  \pause
  O vetor $\vec{u}$ partindo do ponto $(-1, -2)$ termina no ponto com coordenadas
  \[
    \pause
    x
    \pause
    \;=\; -1 + u_x
    \pause
    \;=\; -1 + 3
    \pause
    \;=\; 2
  \]
  \[
    \pause
    y
    \pause
    \;=\; -2 + u_y
    \pause
    \;=\; -2 + 4
    \pause
    \;=\; 2
  \]
\end{resolution}

%------------------------------------------------------------------------------%
\endinput

%------------------------------------------------------------------------------%
\begin{question}
  Avalie e esboce o vetor
  \[
    v(x, y) = \vetor{2x+y-1}{1+y-x}
  \]
  nos pontos
  \(( 0,  0)\),
  \(( 1,  1)\),
  \(( 1, -1)\)
\end{question}

%------------------------------------------------------------------------------%
\begin{resolution}{Ponto \((0, 0)\)}
\begin{columns}
\column{0.5\linewidth}
  \begin{align*}
    \onslide<+->{v(0, 0)}
    \onslide<+->{& = \vetor{2x+y-1}{1+y-x} \evalat{(0, 0)}{} \\}
    \onslide<+->{& = \vetor{2\times 0+0-1}{1+0-0} \\}
    \onslide<+->{& = \vetor{-1}{1}}
  \end{align*}
\column{0.5\linewidth}
  \onslide<+->{\includegraphics{E_01_exemplo_4_vetor_1}}
\end{columns}
\end{resolution}

%------------------------------------------------------------------------------%
\begin{resolution}{Ponto \((1, 1)\)}
\begin{columns}
\column{0.5\linewidth}
  \begin{align*}
    \onslide<+->{v(1, 1)}
    \onslide<+->{& = \vetor{2x+y-1}{1+y-x} \evalat{(1, 1)}{} \\}
    \onslide<+->{& = \vetor{2\times 1+1-1}{1+1-1} \\}
    \onslide<+->{& = \vetor{2}{1}}
  \end{align*}
\column{0.5\linewidth}
  \onslide<+->{\includegraphics{E_01_exemplo_4_vetor_2}}
\end{columns}
\end{resolution}

%------------------------------------------------------------------------------%
\begin{resolution}{Ponto \((1, -1)\)}
\begin{columns}
\column{0.5\linewidth}
  \begin{align*}
    \onslide<+->{v(1, -1)}
    \onslide<+->{& = \vetor{2x+y-1}{1+y-x} \evalat{(1, -1)}{} \\}
    \onslide<+->{& = \vetor{2\times 1+(-1)-1}{1+(-1)-1}  \\}
    \onslide<+->{& = \vetor{0}{-1}}
  \end{align*}
\column{0.5\linewidth}
  \onslide<+->{\includegraphics{E_01_exemplo_4_vetor_3}}
\end{columns}
\end{resolution}

%------------------------------------------------------------------------------%


%------------------------------------------------------------------------------%
%\section{Lista Mínima}
%
%%------------------------------------------------------------------------------%
%\begin{frame}{Lista Mínima}
%  \begin{enumerate}
%    \item
%  \end{enumerate}
%
%  \vfill
%  Atenção: A prova é baseada no livro, não nas apresentações
%\end{frame}

%------------------------------------------------------------------------------%
\end{document}
%------------------------------------------------------------------------------%


